\documentclass[10pt,twocolumn]{article}

\usepackage[margin=0.85in]{geometry}
\usepackage{times}
\usepackage{booktabs}
\usepackage{amsmath}
\usepackage{listings}
\usepackage{xcolor}
\usepackage{url}
\usepackage[hidelinks]{hyperref}
\usepackage{caption}

\captionsetup{font=small,labelfont=bf}

\lstset{
  basicstyle=\ttfamily\scriptsize,
  columns=fullflexible,
  breaklines=true,
  frame=single,
  framesep=3pt,
  rulecolor=\color{gray!50},
  backgroundcolor=\color{gray!5},
  xleftmargin=2pt,
  aboveskip=6pt,
  belowskip=6pt,
}

\title{\vspace{-2em}\textbf{Tritium: One Weight Format for Ternary Inference\\on Commodity CPUs and Multiplier-Free Hardware}}

\author{
Bahadir Kosapinar\\
\small\texttt{https://github.com/jackthepunished/tritium}
}

\date{\small August 2026}

\begin{document}
\maketitle

\begin{abstract}
\noindent
BitNet b1.58 models constrain every weight to $\{-1, 0, +1\}$, which replaces
multiply-accumulate with select-accumulate and removes the multiplier from the
dominant operation of transformer inference. Existing runtimes realise this on
hardware that has multipliers anyway, using weight layouts designed for a CPU.
We describe \emph{Tritium}, a ternary inference stack built around a single
weight layout, \texttt{.trit} v1, in which one 16-byte \emph{beat} is
simultaneously one iteration of a SIMD kernel's inner loop and one cycle's
weight input to a multiplier-free hardware core. A memory-mapped model file is
therefore not decoded into a runtime representation; it is the representation,
on both targets.

We report the format, the CPU kernels, the SystemVerilog core, and a
verification methodology in which three independent implementations of the same
arithmetic are required to agree \emph{exactly} rather than within a tolerance.
On BitNet b1.58 2B4T the runtime reproduces HuggingFace reference logits at mean
cosine 0.9991 with 100\% top-1 agreement, the reference and production
implementations agree at cosine 1.000000 at every position, and the hardware
core produces byte-identical output to the CPU path while Yosys asserts zero
multiplier cells on every build.

On an AMD Ryzen 9 8945HX it decodes at 27.9 tokens/s, ahead of both llama.cpp
and bitnet.cpp up to four threads and 1.43$\times$ bitnet.cpp per core, and
behind bitnet.cpp above four threads where our curve peaks and declines. We
report that asymmetry, and the measurement errors we made and corrected while
producing it, because we think the interesting contribution is not throughput
but the co-designed layout and the verification discipline that makes a hardware
path credible before fabrication.
\end{abstract}

\section{Introduction}

Autoregressive decoding at batch size 1 is memory-bound. Throughput is bounded
by
\begin{equation}
\text{tok/s} \approx \frac{\text{effective memory bandwidth}}{\text{bytes touched per token}}
\end{equation}
because every weight is read once per token and almost no arithmetic is reused.
The standard response is to shrink the weights, which is what four-bit and
eight-bit post-training quantisation do.

Ternary models attack a different term. BitNet b1.58~\cite{bitnet158} trains
with weights constrained to $\{-1, 0, +1\}$, so the inner product
\begin{equation}
y_r = \sum_c W_{rc}\, x_c, \qquad W_{rc} \in \{-1,0,+1\}
\end{equation}
contains no multiplication at all:
\begin{equation}
y_r = \sum_{c\,:\,W_{rc}=+1} x_c \;-\; \sum_{c\,:\,W_{rc}=-1} x_c .
\end{equation}
Each weight participates additively, subtractively, or not at all. We refer to
this as \emph{select-accumulate}. On the 2B4T checkpoint 42.19\% of weights are
exactly zero, so nearly half of the terms are skips.

This is well known. What is less examined is that the runtimes which execute
these models store weights in layouts chosen for a CPU, and a datapath with no
multipliers in it would have to rewrite those bytes before it could consume
them. The arithmetic is hardware-friendly; the encoding is not.

Tritium~\cite{tritium} is an attempt to remove that seam. Its contributions
are:

\begin{enumerate}
\item \textbf{A bit-plane weight layout} (Section~\ref{sec:format}) whose
      16-byte unit is simultaneously one SIMD inner-loop iteration and one
      hardware cycle's weight input, at exactly 2.000 bits per weight.
\item \textbf{A multiplier-free hardware core} (Section~\ref{sec:rtl}) that
      consumes that layout directly from a memory-mapped file, verified
      bit-exact against a software reference and asserted free of multiplier
      cells by synthesis on every build.
\item \textbf{An exactness-based verification methodology}
      (Section~\ref{sec:verification}) in which three independent
      implementations must produce identical integers. We give a concrete
      defect this caught that no tolerance-based test would have.
\item \textbf{An evaluation that includes its own negative results}
      (Section~\ref{sec:eval}), including the thread counts at which Tritium
      loses, and three measurement errors we made and corrected while producing
      it.
\end{enumerate}

\section{Background}

\subsection{The model}

All measurements use \texttt{microsoft/bitnet-b1.58-2B-4T}, a 2.4B-parameter
decoder-only transformer: 30 layers, hidden size 2560, FFN width 6912, 20
attention heads with 5 key/value heads (4:1 grouped-query attention), vocabulary
128256, squared-ReLU activation, tied input and output embeddings. Ternary
weights appear in all seven projections per layer; norms, embeddings and the
language-model head are dense.

\subsection{Where the bytes are}

Table~\ref{tab:bytes} gives the per-token traffic. The result is not what a
ternary-focused design anticipates: even in bf16 the dense tied head rivals
every ternary projection combined, and it exceeded them 2.52:1 while we stored
it widened to f32. Any runtime that vectorises the ternary kernel and leaves the
head alone has not moved end-to-end throughput.

\begin{table}[t]
\centering
\small
\begin{tabular}{lrr}
\toprule
 & bytes/token & share \\
\midrule
Ternary weights (30 layers)     & 521{,}011{,}200   & 44\% \\
Tied head / embeddings (bf16)   & 656{,}670{,}720   & 56\% \\
\midrule
\textbf{Total}                  & \textbf{1{,}177{,}681{,}920} & \\
\midrule
KV cache (f32, 2048 ctx)        & 314{,}572{,}800   & resident \\
\bottomrule
\end{tabular}
\caption{Per-token weight traffic, BitNet b1.58 2B4T.}
\label{tab:bytes}
\end{table}

\section{The \texttt{.trit} v1 layout}
\label{sec:format}

A ternary weight has three states, so two bits is the natural budget under any
encoding. The question is how those two bits are arranged.

\subsection{Interleaved codes versus bit planes}

The obvious encoding assigns a 2-bit code per weight (\texttt{00}$=0$,
\texttt{01}$=+1$, \texttt{10}$=-1$), packing four weights per byte. Extracting
the sign of lane $k$ then requires a shift by $2k$ and a mask, per weight. This
defeats vectorisation: there is no SIMD instruction whose operand is
``every other pair of bits, sign-extended''. An earlier version of this system
used that encoding and ran a branch per weight at 0.16 tokens/s.

Tritium instead stores weights as \emph{bit planes}. A \emph{beat} is 16 bytes
carrying 64 consecutive columns of one row as two 64-bit masks:

\begin{lstlisting}
bytes 0..8    pos   bit k set => W[r][c+k] = +1
bytes 8..16   neg   bit k set => W[r][c+k] = -1
\end{lstlisting}

Neither bit set encodes zero. Both bits set is \emph{invalid} and is rejected at
load and flagged by a sticky error in hardware. Defining it as an error rather
than assigning it a value is deliberate: a kernel that accumulates the two
planes separately would compute $+x-x=0$, while a lane that tests \texttt{pos}
first would compute $+x$. Making it illegal is what prevents the software and
hardware paths from disagreeing on a file neither should have accepted.

Under this layout each of the two sums in Equation~3 is a \emph{mask applied to
a vector of activations}. On AVX-512 the mask is literally a \texttt{k}
register, and the operation is two masked loads and two \texttt{vpdpbusd}. The
decision ``does this weight participate, and with which sign'' costs no branch,
no table lookup, and no per-weight arithmetic.

\subsection{The property that matters}

The bit count is identical to the interleaved encoding. What changes is that
\textbf{one beat is simultaneously one iteration of the CPU kernel's inner loop
and one cycle's weight input to the hardware core}. A memory-mapped
\texttt{.trit} file is not a container that the runtime decodes into some other
layout; it \emph{is} the layout. Nothing between the page cache and the
accumulator rewrites a byte, on either target.

This is the design commitment the rest of the system is organised around. The
CPU and the hardware are not two implementations of one idea kept in
correspondence by discipline. They consume the same bytes.

\subsection{Invariants}

Two invariants are checked on demand and are load-bearing for both targets.

\textbf{P1: planes are disjoint.} $\texttt{pos} \wedge \texttt{neg} = 0$ for
every beat.

\textbf{P2: padding is clear.} When \texttt{cols} is not a multiple of 64, the
final beat of each row covers columns that do not exist. Those bits must be zero
in both planes. With
\begin{equation*}
\texttt{tail} =
\begin{cases}
\texttt{UINT64\_MAX} & \texttt{cols} \bmod 64 = 0\\
(\texttt{UINT64\_C(1)} \ll (\texttt{cols} \bmod 64)) - 1 & \text{otherwise}
\end{cases}
\end{equation*}
every row's last beat satisfies
$(\texttt{pos} \mid \texttt{neg}) \wedge \sim\!\texttt{tail} = 0$.

P2 is what makes zero padding \emph{exact} rather than conventional: a padded
column contributes nothing, so a kernel may process whole 64-column blocks
unconditionally and never needs a scalar tail loop. Both invariants are cheap to
establish at write time and $O(\text{payload})$ to verify.

At column counts that are multiples of 64, which every projection in this model
class satisfies, the layout costs exactly 2.000 bits per weight.

\section{CPU kernels}

Kernels exist for AVX-512 VNNI, AVX-512BW, AVX2, NEON, NEON with the
\texttt{dotprod} extension, and a portable scalar path. All read the planes in
place. Because the accumulators are \texttt{i32} and integer addition is exact
and order-independent, every kernel is \emph{bit-identical} to the portable one.
That is what permits the differential tests of
Section~\ref{sec:verification} to demand equality rather than a tolerance, and
why changing kernels can never move a logit.

Kernel selection is resolved once by runtime feature detection. An override
pins a kernel by name and \emph{errors} if the CPU cannot run it rather than
falling back, so a benchmark can never silently measure a path other than the
one it names, and a CI runner without AVX-512 fails loudly instead of
re-testing the scalar path under a different label.

\begin{table}[t]
\centering
\small
\begin{tabular}{lrrr}
\toprule
Kernel & ms & GB/s & vs scalar \\
\midrule
scalar        & 9.10 & 0.5  & 1.00$\times$ \\
avx2          & 0.66 & 6.7  & 13.87$\times$ \\
avx512bw      & 0.38 & 11.6 & 23.91$\times$ \\
avx512vnni    & 0.28 & 15.8 & \textbf{32.45}$\times$ \\
bitserial     & 6.12 & 0.7  & 1.49$\times$ \\
\bottomrule
\end{tabular}
\caption{Ternary matvec on a $2560\times6912$ tensor, planes resident in L3.}
\label{tab:kernels}
\end{table}

\subsection{On popcount}

Population count is the natural primitive when \emph{activations} are also
1-bit. Here weights are 1.58-bit but activations are int8, so a popcount of a
weight mask only counts participating terms and cannot recover their sum. The
production kernels mask activation bytes and accumulate the two planes
separately.

The bit-serial popcount formulation is nonetheless implemented, because it is
the formulation that maps directly onto the hardware adder tree. It measures
22$\times$ slower than the AVX-512 path at int8 activations
(Table~\ref{tab:kernels}). It is retained as a third independent implementation
for differential testing, and because it becomes the correct kernel if
activations ever drop below 8 bits.

\section{The hardware core}
\label{sec:rtl}

\texttt{trit\_matvec.sv} is a streaming ternary matrix-vector core. Weights
arrive as two 64-bit planes per beat, matching the \texttt{.trit} v1 payload
byte for byte. Activations are preloaded as int8. Each of 64 lanes muxes
$\{+x, -x, 0\}$ into a combinational adder tree feeding a 32-bit accumulator.

Two details are worth stating because they are where naive implementations
diverge from software. First, lanes sign-extend to the accumulator width
\emph{before} negating: $x = -128$ has no positive counterpart in eight bits,
and any implementation that negates in the int8 domain produces a different
answer on that input. The golden vector set includes an \texttt{extremes} case
that feeds $-128$ deliberately, and it must never be weakened. Second, the
four lane states are expressed as a \texttt{unique case} rather than a nested
conditional, which tells synthesis they are mutually exclusive; written as an
if-chain the module synthesises to 48.4k cells instead of 33.7k.

Yosys generic synthesis reports 33{,}659 cells at 64 lanes, of which
approximately 4.1k flip-flops are the flattened activation memory (block RAM on
any real device). \textbf{The build asserts zero \texttt{\$mul} and zero
\texttt{\$macc} cells on every run.} Norm folding
(Section~\ref{sec:numerics}) additionally removes the reciprocal square root
and the divide from the element datapath.

\subsection{Hardware in the loop}

The core is exposed to the runtime as an ordinary backend behind the same trait
as the CPU kernels, driven through a C shim over the Verilated model. Because
v1 stores planes in the core's beat order, mapped model bytes stream in
directly with no repacking between the file and the device under test.

The full 2B4T model decodes end to end through this path, producing
\textbf{byte-identical output} to the CPU backend, at 24.4 seconds per token
under simulation. That figure is a simulator speed, not a device speed, and no
FPGA performance is claimed anywhere in this paper.

\section{Numerics}
\label{sec:numerics}

Three evaluation modes form a ladder, each a strictly more exact evaluation of
the same model, selected automatically to the best the architecture supports.

\textbf{Reference.} Textbook. RMSNorm materialised, activations quantised per
matvec.

\textbf{Folded.} The per-element $1/\text{rms}$ divide leaves the datapath
entirely. Absmax quantisation codes are invariant under a uniform positive
scale, so the codes may be computed from $x \odot g$ directly and the rms
survives as a single per-token scalar folded into the output scale. This is a
numerics result and a synthesis result at once: it is why the hardware needs no
divider.

\textbf{IntMlp.} The squared-ReLU stage never exists in floating point. With
$g = a_g S_g$ and $u = a_u S_u$ from the gate and up projections and uniform
positive scales, $\mathrm{relu}(g)^2 u = t K$ for integer
$t = \mathrm{relu}(a_g)^2 a_u$ with $|t| < 2^{55}$, exact in \texttt{i64}, and
uniform $K = S_g^2 S_u$.

Both upper rungs were validated against the real checkpoint before being made
the default. Measured activation ranges motivated the design: the residual
stream reaches $\sim$138k and the squared-ReLU stage $\sim$1.5$\times10^{10}$,
so int16 is not viable for those stages and the wide accumulators are not
conservatism.

\section{Verification}
\label{sec:verification}

Three independent implementations of the same arithmetic must agree:

\begin{itemize}
\item \textbf{The oracle.} A deliberately naive scalar implementation that
      expands bit planes to one \texttt{i8} per weight. Its job is to be
      evidently correct, not fast.
\item \textbf{The runtime.} Zero-copy, vectorised, threaded.
\item \textbf{The hardware core}, under Verilator.
\end{itemize}

Agreement on the integer path is demanded \emph{exactly}. Lane order, thread
count and the hardware's beat-serial accumulation must all produce identical
bits. This is achievable by construction rather than by tuning, because integer
addition is exact and order-independent; an implementation that cannot meet it
is wrong, not approximate. Only the floating-point tail (attention, the head)
admits reassociation.

\subsection{Why exactness, concretely}

Tolerance-based testing fails on this workload because the failure mode is
silent and delayed. A wrong logit does not crash; it perturbs a probability, and
the divergence appears many tokens later as a different word.

The runtime once folded two scale factors into a single constant and multiplied
once, where the reference multiplies by each in turn. Floating-point
multiplication is not associative, so $a(bc)$ and $(ab)c$ round differently. The
observed consequence: cosine similarity 1.000000 for four token positions,
0.999624 by position seven, and at token sixteen the generated text read
``known'' where the reference read ``the''. No tolerance that admits legitimate
reassociation in the attention tail would have flagged the first four positions,
and no human reading either the code or the output would have found it.

\subsection{Gates}

Results are recorded in a frozen gate document. Correctness gates are
invariants: they may not move at all. Resource gates are targets and are
expected to move, but only accompanied by a written argument naming the change
and its before/after. Table~\ref{tab:gates} lists the measured values.

\begin{table}[t]
\centering
\small
\begin{tabular}{ll}
\toprule
Gate & Measured \\
\midrule
Logits vs HuggingFace     & cosine 0.9991, top-1 100\% \\
Oracle vs runtime         & cosine 1.000000, all positions \\
Greedy text               & byte-identical, all modes \\
Hardware vs golden vectors& bit-exact, incl.\ $x=-128$ \\
Synthesis                 & 0 \texttt{\$mul}, 0 \texttt{\$macc} \\
Hardware in the loop      & byte-identical to CPU \\
\bottomrule
\end{tabular}
\caption{Correctness gates. All are invariants.}
\label{tab:gates}
\end{table}

\subsection{Defects this methodology found}

Beyond the scale-folding bug, two are worth reporting because both were latent
in code that a test matrix claimed to cover.

The NEON kernel accumulated into \texttt{i16} and flushed every 64 \emph{beats},
under a comment reasoning about 64 \emph{accumulation steps}. Each beat issues
four pairwise-add-accumulate instructions, so the interval was 256 steps against
a safe bound of 128; the per-step bound was also understated, since two lanes of
$-128$ pairwise-add to $-256$. The worst case, a 6912-column row of $-1$ weights
against $-128$ activations, wrapped silently and produced a wrong row. It had
never executed: the aarch64 CI job had failed at a linter since it was added and
never reached an assertion, while the matrix advertised ARM coverage. A job that
cannot reach its assertions is not coverage.

Separately, streaming detokenisation routed each token through a decode call
that terminates in a lossy UTF-8 conversion. On this vocabulary U+1F600 splits
across token 76460 (bytes \texttt{f0 9f 98}) and token 222 (byte \texttt{80});
decoding each alone returns the replacement character, so every emoji, flag and
multi-byte sequence emitted by the streaming API was corrupted while a one-shot
decode of the same sequence was correct.

\section{Evaluation}
\label{sec:eval}

\subsection{Method}

All measurements are on an AMD Ryzen 9 8945HX (Zen 4, AVX-512 VNNI, 32 threads,
DDR5), greedy decoding, 32 tokens, four-prompt suite, page cache pre-warmed,
median of three invocations. Raw data is committed to the repository.

Three methodological points, each of which changed a published number:

\textbf{Median, not best.} Reporting our best run against a baseline's median
tilts every comparison. All runtimes are now reported as medians. Where a
speedup is claimed it is an interleaved paired measurement, because this host
drifts by up to 30\% between runs and sequential blocks cannot survive that.

\textbf{Pre-warm.} The model file lives on a mount whose cold reads are slow. A
cold measurement read 11.25 tokens/s where three consecutive warm ones read
13.5--13.9. An earlier version of Table~\ref{tab:compare} claimed Tritium was
fastest single-threaded; that was an artifact of our model being warm in page
cache while freshly downloaded baseline models were not.

\textbf{A measured roofline.} Achieved bandwidth is divided by a streaming-read
probe run on the same machine in the same invocation, best of five passes. A
single-pass probe under-reported by 10--15\%, which inflated every reported
fraction. The probe is still noisy after the fix, reading 48--58 GB/s across the
runs in one sweep, so we quote it alongside the fraction it produced rather than
claiming a fixed ceiling for the host.

\subsection{Absolute performance}

\begin{table}[t]
\centering
\small
\begin{tabular}{lr}
\toprule
Decode                    & 27.90 tok/s \\
Time to first token       & 239 ms \\
Peak RSS                  & 1220 MB \\
Achieved bandwidth        & 32.9 GB/s \\
Probe, same invocation    & 57.2 GB/s \\
Fraction of it            & 57\% \\
\bottomrule
\end{tabular}
\caption{Tritium, 4 threads.}
\label{tab:abs}
\end{table}

For reference, the pre-format-change system decoded at 0.16 tokens/s with a peak
resident set of 6.10 GiB.

\subsection{Comparison}

\begin{table}[t]
\centering
\small
\begin{tabular}{rrrr}
\toprule
threads & Tritium & llama.cpp & bitnet.cpp \\
        & ternary & Q4\_K\_M  & I2\_S \\
\midrule
1  & \textbf{19.40} & 15.40 & 13.61 \\
2  & \textbf{26.73} & 22.06 & 19.10 \\
4  & \textbf{27.90} & 27.83 & 26.72 \\
8  & 25.93 & 26.18 & \textbf{32.49} \\
16 & 20.72 & 23.88 & \textbf{31.21} \\
\bottomrule
\end{tabular}
\caption{Decode tokens/s. bitnet.cpp runs the identical checkpoint. llama.cpp
runs Qwen2.5-3B Q4\_K\_M, since mainline cannot load the \texttt{i2\_s} type and
its converter has no BitNet entry: a different model at a different quality
point, included for context rather than as a like-for-like comparison.}
\label{tab:compare}
\end{table}

Table~\ref{tab:compare} splits cleanly. Tritium leads at one, two and four
threads and is 1.43$\times$ bitnet.cpp per core on the identical checkpoint.
Above four threads it loses: its curve peaks and declines while bitnet.cpp keeps
climbing to 32.49. Scaling from one thread to eight is 1.34$\times$ against
bitnet.cpp's 2.39$\times$.

An earlier version of this table had Tritium last per core. Two changes moved
it: a persistent worker pool worth 1.17$\times$, and storing the head in the
bf16 the checkpoint already used rather than widened f32, worth 1.37$\times$.
The second was a defect rather than an optimisation. bitnet.cpp already stored
its embedding table in f16, and that was most of its remaining advantage.

\subsection{Diagnosis: bytes are not time}

An instrumented decode contradicts the byte analysis. Measured before either
change in Section~\ref{sec:eval} landed, with an f32 head and a
single-threaded ternary path: the head was 42\% of wall time despite being
72\% of the bytes at that point, because the two halves ran at opposite ends of
the machine's efficiency range, 48.8 GB/s against 13.9 GB/s. Both figures have
since moved -- the head is bf16 and 56\% of the bytes, and the ternary path
uses the pool -- but the lesson is why they were worth measuring separately.

\begin{table}[t]
\centering
\small
\begin{tabular}{lrrr}
\toprule
 & per token & share & achieved \\
\midrule
30 transformer layers & 37.5 ms & 58\% & 13.9 GB/s \\
Language-model head   & 26.9 ms & 42\% & 48.8 GB/s \\
\bottomrule
\end{tabular}
\caption{Where decode time was spent at 4 threads, \emph{before} the bf16 head
and the persistent pool. Retained because both conclusions drawn from it held.}
\label{tab:time}
\end{table}

Two opposite conclusions followed, and both held. The head was close enough to
what the memory system delivers that no kernel work moved it: forcing the dense
kernel to scalar, AVX2 and AVX-512 in turn gave 28.8, 28.1 and 27.1 ms. Only
moving fewer bytes helped, which is what storing it in bf16 does. The ternary
path was the opposite, compute-bound in the kernel at approximately 15.8 GB/s
per core and barely above what it achieved from DRAM, so it would scale with
cores if the runtime could use them.

\subsection{Why it does not scale}

A decode step issues 210 ternary matvecs whose largest is 4.4 MB, and each one
paid a work-stealing fork/join. Lowering the threshold at which matvecs
parallelise made matters dramatically worse, not better
(Table~\ref{tab:par}, row \emph{rayon-global}): even two threads was
1.8$\times$ slower than not parallelising at all.

\begin{table}[t]
\centering
\small
\begin{tabular}{lrrrr}
\toprule
Mechanism & 1t & 2t & 4t & 8t \\
\midrule
serial (previous) & 37.6 & 37.6 & 36.9 & 36.7 \\
rayon, global pool& 36.0 & 65.0 & 108.8& 174.0 \\
rayon, sized pool & 35.8 & 50.0 & 69.7 & 148.9 \\
rayon, broadcast  & 35.9 & 47.2 & 57.5 & 89.8 \\
persistent spin barrier & 37.1 & \textbf{20.0} & \textbf{12.8} & 14.1 \\
\bottomrule
\end{tabular}
\caption{Layer time in ms by parallel mechanism. All produce byte-identical
output.}
\label{tab:par}
\end{table}

Table~\ref{tab:par} ruled out the cheap explanations. A dedicated pool sized to
the requested thread count helped substantially and was still worse than serial
execution. A broadcast primitive helped more and was still worse. Only a
persistent worker pool on a sequence counter, the mechanism ggml uses, beat
serial execution, at 2.9$\times$ on the layer path.

That pool is what ships. Measured end to end against the previous
implementation in interleaved pairs, because this host drifts by up to 30\%
between runs, it is worth \textbf{1.17}$\times$ and lifts one-to-eight scaling
from 1.15$\times$ to 1.34$\times$. An isolated prototype had suggested
1.59$\times$; the smaller figure is the one measured on the whole decode, and it
is the one we report.

\section{Limitations}

We state these because several are the first thing a reader should check.

\begin{itemize}
\item \textbf{Throughput above four threads.} Decode peaks at four threads and
      declines, where bitnet.cpp climbs to eight and reaches 32.49 against our
      25.93. The persistent pool improved one-to-eight scaling from
      1.15$\times$ to 1.34$\times$ and did not close the gap.
\item \textbf{No energy measurement.} Joules per token is the metric this
      architecture exists to improve and we have never reported one, because no
      machine in our loop exposes a counter. The field is left empty rather than
      estimated.
\item \textbf{No silicon.} The hardware core is simulation-first and not
      timing-closed. The single-cycle 64-term reduction and 64 parallel
      activation reads are fine under Verilator and unproven on a device.
      Activation memory is capped at 8192 columns.
\item \textbf{No ARM timing.} The NEON kernels are verified bit-exact on
      aarch64 hardware in CI and have never been timed. No ARM performance
      figure exists.
\item \textbf{Scope.} Batch size 1, no batched prefill, context capped at 2048
      with an f32 preallocated KV cache.
\end{itemize}

\section{Related work}

\texttt{bitnet.cpp}~\cite{bitnetcpp} is the reference ternary runtime and the
comparison point in Table~\ref{tab:compare}. It is faster than Tritium on CPU
today. It targets CPUs, and its \texttt{i2\_s} tensor type is specific to that
fork; mainline \texttt{llama.cpp}~\cite{llamacpp} cannot load it. Ternary tensor
types exist in mainline (\texttt{TQ1\_0}, \texttt{TQ2\_0}) but its HuggingFace
converter has no BitNet entry.

The distinction we draw is not throughput but target. These are fast CPU
runtimes whose weight layouts serve a CPU. Tritium's layout is a hardware
interface that a CPU also happens to consume efficiently, which is what makes
the path in Section~\ref{sec:rtl} available without a repacking step.

\section{Future work}

An energy measurement on a machine with a power counter is first, because it
converts the central claim of this architecture from an argument into a number,
and because nothing else on the list changes what the project is. Then the
decline above four threads, which the persistent pool improved without
eliminating. Then ARM throughput, where the kernels are verified correct and
have never been timed. Then board bring-up.

\section{Conclusion}

Ternary weights make the multiplier optional, but only if the weights are stored
in a form a multiplier-free datapath can consume. Tritium's contribution is a
layout in which one 16-byte unit is both a SIMD inner-loop iteration and a
hardware cycle, so a memory-mapped model file is the interface to both targets
rather than an input to be repacked for either.

We do not claim throughput leadership; on CPU we are measurably behind, for a
diagnosed and prototyped reason. We claim that the arithmetic is proven correct
against an independent reference, that the hardware core is proven bit-exact
against that same reference and free of multipliers by construction, and that
this is the order in which a silicon architecture should be de-risked: the
expensive irreversible step last, not first.

\begin{thebibliography}{9}
\small
\bibitem{bitnet158}
S.~Ma et al.,
``The Era of 1-bit LLMs: All Large Language Models are in 1.58 Bits,''
arXiv:2402.17764, 2024.

\bibitem{bitnetcpp}
Microsoft,
``BitNet: Inference framework for 1-bit LLMs,''
\url{https://github.com/microsoft/BitNet}.

\bibitem{llamacpp}
G.~Gerganov et al.,
``llama.cpp,''
\url{https://github.com/ggml-org/llama.cpp}.

\bibitem{tritium}
B.~Kosapinar,
``Tritium,''
\url{https://github.com/jackthepunished/tritium}.
\end{thebibliography}

\end{document}
